\documentclass[10pt]{article}
\usepackage[margin=20mm]{geometry}
\usepackage{amsmath,amssymb,mathtools,fontspec,unicode-math,xurl,hyperref}
\IfFontExistsTF{DejaVu Serif}{\setmainfont{DejaVu Serif}}{\setmainfont{Latin Modern Roman}}
\setmathfont{latinmodern-math.otf}
\setlength{\parindent}{0pt}\setlength{\parskip}{6pt}\sloppy
\title{Initial scientific review of
Poisson spacing statistics for lattice points on circles}
\author{}
\date{26 October 2026}
\begin{document}
\maketitle
Abstract: Two explicit proof steps require correction. The reindexing in Proposition 5.8 replaces a constrained stratum by a different, potentially nonempty stratum and counts some leading contributions twice. The final density argument applies a quoted multiplicative-function estimate to a function outside the stated hypotheses. Exact low-dimensional checks and analytic calculations establish both discrepancies. Retaining only the original strata that satisfy the constraint repairs the first step; counting distinct split prime factors repairs the second. Both steps feed the density-one spacing conclusion and are rated Major under the specified core-derivation standard. Neither main theorem statement is refuted. A separate complex-exponent scope question is left unrated. Independent scientific acceptance remains pending.\par
\section*{Identity, version and coverage}
The assigned paper is Pär Kurlberg and Stephen Lester, Poisson spacing statistics for lattice points on circles, M25-0064, Inventiones mathematicae 240 (2025), 537–585, DOI 10.1007/s00222-025-01324-1 [paper]. The reviewed original is the acquired publisher version of record, 49 PDF pages. Throughout this report, printed page \(p\) corresponds to one-based PDF page \(p-536\). Its SHA-256 is\par
All 3,029 lines of its native text extraction, including references, were read continuously. Nineteen original PDF pages were inspected visually, including every page containing a figure (PDF pages 3 and 24), the theorem statements, both findings, and the concluding proof. This is claim-focused reading and checking, not an independent reconstruction of every external analytic theorem.\par
The acquired arXiv:2112.08522v1 source, dated 15 December 2021, is an older author version [arxiv]. Targeted static inspection of its `merge-arxiv.tex`, lines 1071–1096 and 2167–2200, corroborated the complex-exponent wording and the reindexing expression. It was not treated as the source of the journal PDF, compiled, or executed. All 12 manifest-listed files had actual byte counts and SHA-256 values checked. No separate essential supplement is missing. The absence of an identified paper-specific program is not a defect.\par
\section*{Claims, hypotheses and dependencies}
Theorems 1.1 and 1.2 assert exponential nearest-neighbor and joint consecutive-gap laws along a density-one subsequence of the sums of two squares. Density is relative to \(\mathcal S\), and the number \(\ell\) of consecutive gaps is fixed. The angle variable is modulo \(\pi/2\), \(N=r(n)=r_2(n)/4\), and the gap scale is \(N/(\pi/2)\). The result does not assert Poisson behavior along every sequence of radii or uniformly for growing \(\ell\).\par
The proof freezes a square factor and a fixed integer \(n_0\), writing \(n=n_0n_1\) with \(n_1\) square free, odd, composed of split primes, and coprime to \(n_0\). If \(M\) is its number of prime factors, then \(N=2^M r(n_0)\). Theorem 3.1 matches the first two arithmetic moments to those of a random subset-sum model, with error \[O_{f,r,n_0,A}\bigl((\log\log x)^{-1}+M^{-10}\bigr),
 \qquad 1\le M\le A\log\log x.\] The constants do not assert uniformity for arbitrary growing \(n_0\). The final slowly growing cutoff and diagonalization are needed for that reason.\par
Proposition 4.1 supplies the arithmetic average. Propositions 5.1 and 5.8 remove its arithmetic factor after averaging frequencies. Section 6 identifies the resulting expressions with the random model. The rank estimates in Section 7 give Poisson correlations; Section 8 combines variance, density estimates, sieving, and diagonalization to obtain spacings. Problem 1 lies on the second-moment branch of this chain. Problem 2 concerns the final density-one selection itself.\par
\section*{Problem 1: adding a relation does not restrict a stratum}
\subsection*{Source and suspect step}
In Section 5.1, printed p. 559 (PDF 23), the paper defines \[V_S^\star=\ker M_S\setminus
       \bigcup_{S_1\supsetneq S}\ker M_{S_1}.\] For a nonempty stratum, \(S\) is the exact set of complementary subset classes whose associated forms vanish. Thus forms outside \(S\) are required to be nonzero; this is stronger than merely belonging to \(\ker M_S\).\par
In the proof of Proposition 5.8, equation (5.36), p. 569 (PDF 33), the sum over each \(V_S^\star\) still has the constraint \(k_r=0\). On p. 570 (PDF 34), first paragraph and first unnumbered display, the paper sets \[\mu=[\{0,\ldots,r-1\}],\qquad
 \widetilde S=S\cup\{\mu\},\qquad
 \widetilde V^\star=V_{\widetilde S}^\star\] and replaces the constrained sum by a sum over \(\widetilde V^\star\), while retaining the outer sum over all original \(V^\star\in\mathcal V_{2r}\). This replacement is false. If \(\mu\notin S\), the original stratum requires \(k_\mu\ne0\), so its intersection with \(k_r=k_\mu=0\) is empty. Adding \(\mu\) to the label instead creates a different stratum, which can be nonempty and already appears elsewhere in the outer sum.\par
\subsection*{Complete countercheck with all parameters}
Take correlation order \(r=2\) and \(n_0=1\). These are allowed by Proposition 5.8; \(M\) can be any positive integer in its stated range, and \(N=2^M\). The frequency space for the second moment is \(\mathbb R^3\), with \(k_0=k_4=0\), and \(\mu=[\{0,1\}]\) gives \(k_\mu=k_2\).\par
Represent each class of subsets of \(\{0,1,2,3\}\) by the member not containing 3. Write \(S_0=\{[\varnothing]\}\) and \(S_1=\{[\varnothing],\mu\}\). The generic stratum \(V_{S_0}^\star\) is nonempty but contains no point with \(k_2=0\). In contrast, \[U:=V_{S_1}^\star
   =\{(a,0,b): a b\ne0,\ a\ne b,\ a\ne-b\}.\] Indeed, the four increments are \((a,-a,b,-b)\); under these exclusions the only zero-sum subset classes are \([\varnothing]\) and \(\mu\). The point \((1,0,3)\) supplies an explicit witness. Thus both \(S_0\) and \(S_1\) label nonempty original strata, but the printed mapping sends both to \(S_1\): \[S_0\cup\{\mu\}=S_1,\qquad S_1\cup\{\mu\}=S_1.\] The original constrained sum counts each point of \(U\) once; the replacement counts it twice. This argument uses precisely the primed/stratified definition, not a different convention for ordinary kernels.\par
The independent exact-arithmetic check in `verification/0064/evidence/reviewer\_checks.json` enumerated all 18 nonempty strata for this arrangement. It confirmed the two preimages, the zero classes of \((1,0,3)\), and all cube-intersection bounds \(\alpha_V\le 2^{3-\dim V}\) in this finite case. In the box \(|a|,|b|\le4\), \(U\) has 48 integer points; each has the stated duplicate. These finite calculations corroborate the general set-theoretic argument.\par
\subsection*{The duplicate is of main-term size}
This discrepancy is not removed by deleting lower-dimensional boundaries. Choose the permitted Schwartz test function \(f(t)=e^{-\pi t^2}\), so \(\widehat f=f\), and use \[h(x_1,x_2,x_3)=e^{-\pi(x_1^2+x_2^2+x_3^2)}.\] It satisfies (5.33), \(h(x_1,0,x_3)=\widehat f(x_1)\widehat f(x_3)\). Also \(\widetilde\ell_{\vec k}(1)=1\). On \(U\), \(\alpha_U=2\). After the stratum averaging asserted in the next display of the proof, the extra main contribution equals \[\frac{4^M}{N^4}
 \sum_{\substack{a,b\in\mathbb Z\\ab\ne0,\ a\ne\pm b}}
           e^{-\pi(a^2+b^2)/N^2}
 =
 \frac1{N^2}
 \sum_{\substack{a,b\in\mathbb Z\\ab\ne0,\ a\ne\pm b}}
           e^{-\pi(a^2+b^2)/N^2}
 \longrightarrow 1.\] The limit follows from the two-dimensional Riemann sum; removing four lines changes it by \(O(N^{-1})\). The integral is \(\int_{\mathbb R^2}e^{-\pi(u^2+v^2)}\,du\,dv=1\). The averaging errors tend to zero as \(M\to\infty\) and \(\log\log x\to\infty\) within the allowed range. The other surviving original strata remain represented, and this positive main contribution cannot be absorbed by the stated \(o(1)\) error. The computation here is analytic; the finite enumeration does not purport to prove the limit.\par
\subsection*{Counterevidence, repair and impact}
The exact Fourier expansion before this reindexing is not contradicted, and the statement of Proposition 5.8 is not disproved. A possible informal reading is that the sum is intended to range only over distinct nonempty resulting strata, with duplicates removed. That convention is absent from the displayed outer index and would be a substantive correction to the equality as written.\par
A direct repair is simpler. For every nonempty original stratum, \[V_S^\star\cap\{k_r=0\}=
 \begin{cases}
   V_S^\star,&\mu\in S,\\
   \varnothing,&\mu\notin S.
 \end{cases}                                                    \tag{R1}\] Keep exactly the first set of strata, each once, and do not change \(S\). On each retained stratum the arguments of Lemmas 5.6 and 5.7 apply as before with \(2r\) in place of \(r\). The product \(\widetilde\ell_{\vec k}(n_0)\) is the same \(2r\)-factor product when \(k_r=0\). If the stratum dimension is \(d\), Lemma 5.3 gives \[(2\alpha_V)^M N^{d-2r}
 \le (2^M)^{2r-d}(2^M r(n_0))^{d-2r}
 =r(n_0)^{d-2r},\] which is bounded for fixed \(n_0\). There are only \(O_r(1)\) strata. Summing their existing errors proves the claimed estimate in Proposition 5.8, conditional on the earlier analytic estimates already used there. This verifies the local repair and its error scaling.\par
The affected second moment feeds Theorem 3.1(3.4), then the variance estimate (8.1), and hence the density-one Poisson conclusion. The demonstrated invalid core derivation therefore supports Major under the specified calibration, even though the local repair is short. No counterexample to Theorems 1.1 or 1.2 is claimed.\par
\section*{Problem 2: the final multiplicative-function invocation fails its hypothesis}
\subsection*{Location and exact failed condition}
On p. 543 (PDF 7), Section 2 explicitly defines the multiplicity-counting function \[\Omega_1(n)=\sum_{\substack{p^a\Vert n\\p\equiv1\pmod4}} a.\] On p. 583 (PDF 47), after (8.2), the proof quotes the bound \[\sum_{n\le x}h(n)\ll\frac{x}{\log x}
       \prod_{p\le x}\left(1+\frac{h(p)}p\right)
 \quad\text{for multiplicative }0\le h(n)\le d(n),
                                                               \tag{R2}\] where \(d(n)\) is the divisor count. It then applies that bound with \(h(n)=b(n)2^{\Omega_1(n)}\), where \(b\) is the indicator of sums of two squares.\par
This chosen \(h\) is nonnegative and multiplicative, but it does not satisfy the displayed upper hypothesis. For \(n=25=5^2\), one has \[b(25)=1,\quad \Omega_1(25)=2,\quad h(25)=4,\quad d(25)=3.\] All conventions agree with the source: \(5\equiv1\pmod4\), multiplicity is two, and 25 is a sum of two squares. More generally, \(h(5^a)=2^a>a+1=d(5^a)\) for every \(a\ge2\), with unbounded ratio. Thus even inserting a fixed implied constant in the pointwise domination would not fix it. The saved exact check records exponents \(1\le a\le8\); the general inequality is elementary.\par
This is a counterexample to the invocation’s required hypothesis, not to the desired summatory estimate or density conclusion. If a stronger external multiplicative-function theorem is intended, its broader hypotheses and verification must replace the restricted form actually stated. Reinterpreting \(\Omega_1\) as a distinct-prime count silently changes the explicit definition in Section 2.\par
\subsection*{Verified repair preserving the density estimate}
Define instead \[\omega_1(n)=\#\{p\mid n:p\equiv1\pmod4\},\qquad
 h_*(n)=b(n)2^{\omega_1(n)}.\] This is nonnegative and multiplicative. At split prime powers its value is 2, at powers of 2 it is 1, and at inert prime powers it is 0 or 1. Consequently \(h_*(p^a)\le a+1\) at every prime power, and \(h_*(n)\le d(n)\) for all \(n\). Its values at primes are the same as those of the function used in the source. Applying (R2), with the same split-prime reciprocal-sum estimate as in the paper, gives \[\sum_{n\le x}h_*(n)\ll x.\]\par
For the paper’s factorization \(n=n_0n_1\), the integer \(n_1\) is square free. Therefore \[\Omega_1(n_1)=\omega_1(n_1)\le\omega_1(n).\] When \(n\ge\sqrt x\) and \(\Omega_1(n_1)\ge\log\log n\), the last inequality implies \(2^{\omega_1(n)}\gg(\log x)^{\log2}\), since \(\log\log n\ge\log\log x-\log2\). Markov’s inequality now gives exactly the required scale: \[\#\{n\in\mathcal S:n\le x,\ \Omega_1(n_1)\ge\log\log n\}
 \ll\sqrt x+\frac{x}{(\log x)^{\log2}}
 =o\!\left(\frac{x}{\sqrt{\log x}}\right).\] Thus the auxiliary function can be changed without weakening the density claim. This correction uses the already available square-free property, and does not alter the global definition of \(\Omega_1\).\par
The error occurs in the final proof that the restricted set of radii has full density, an essential part of both main statements. It is a demonstrated invalid application on that core chain, so it is rated Major under the requested standard. The ease of the verified repair and survival of the claimed bound are explicit counterevidence against any stronger conclusion about the paper’s validity.\par
\section*{Additional checked reasoning and unresolved scope}
The first-moment calculation, the cube-intersection bound, the rank-counting strategy, and the distinctions between all tuples and distinct tuples were read and checked at their main logical transitions. For fixed \(r\), generic binary tuples contain the required coordinate patterns with probability tending to one. For a nontrivial low-rank relation, the positive-pair argument in Proposition 7.5 gives \(o(N^{d+1})\) configurations, rather than only \(o(N^r)\). This stronger scale is what is needed against a contribution of size \(N^{-d}\). No new counterexample to that estimate was established.\par
For Lemma 7.8, distinct points may have the same subset \(J\) but different deterministic shifts. This does not alone invalidate the argument: for fixed \(n_0\) the finitely many distinct shifts have positive minimum separation on the torus, while a compactly supported test requires all points in a contributing tuple to be within \(O_{r,f}(N^{-1})\) of one another. Hence such omitted tuples contribute zero for sufficiently large \(N\). The fixed-\(n_0\) qualification is essential.\par
For the extension from a countable test family to all compactly supported smooth functions, positivity and cutoff functions can supply the needed local mass bound. A bound on an expected functional alone should not be interpreted as an automatic uniform bound on every deterministic radius. These details were considered in interpreting the standard diagonalization, without asserting a separate disproved theorem.\par
\subsection*{Complex powers in Proposition 4.4: scope clarification}
Equation (4.17), p. 550 (PDF 14), prints a restriction on \(|\operatorname{Re}z|\), whereas the coefficient bound just above it, (4.14), assumes \(|z|\le A\). The proof of (4.17) is omitted as similar to that of (4.16); the later application in Lemma 5.6 uses fixed real integer powers \(u\).\par
An interpretation uniform over arbitrary imaginary \(z\) would be false. Fix \(k=1\), let \(y=2^J\) be sufficiently large, and put \(z=iT\). The truncated Dirichlet sum is a polynomial in \(z\) of degree \(J\) with nonzero leading coefficient \[\frac{(-1)^J}{2^J J!}.\] The unique degree-\(J\) contribution at norms at most \(2^J\) is the ideal \((1+i)^J\); any replacement of a norm-2 prime ideal by another prime ideal increases the norm. On the other hand, the paper’s \(L(1,4)>0\) implies \(|L(1,4)^{iT}|=1\). Thus the polynomial grows unboundedly with \(|T|\), although \(\operatorname{Re}z=0\) satisfies the printed range. A uniform error depending only on \(y,k,A,n_0\) cannot hold in this interpretation.\par
The prior bounded-\(z\) context and unspecified parameter dependence for this second formula leave an alternative reading. Accordingly this is retained as an unrated scope clarification, not added to the supported issue count. Explicitly retaining bounded \(|z|\) (in particular the fixed real powers actually used) avoids this countercheck. This review does not claim to have verified a larger growing-complex-\(z\) range.\par
The cited estimates for Hecke \(L\)-functions, coefficient extraction, and the external theorem transferring correlations to gaps were not independently reproved from their primary proofs. All extracted source text was read, but this limitation precludes a whole-paper correctness certificate. The plots were visually checked as illustrations; their underlying data were not independently regenerated. They are not needed for the two analytic findings.\par
\section*{Disposition, versions and remaining work}
The acquired journal PDF and the targeted older TeX passages both contain the reindexing under discussion. This is not evidence of an unpublished correction or author adoption. A limited exact-title/DOI correction search and inspection of the publisher and arXiv records did not identify a specific erratum. The arXiv record exposes v1; no later author file was substituted for the frozen input. The search does not exclude all later commentary or corrections.\par
Item | Judgment | Repair verification | Adoption / remaining work\par
F01 | Major | Correct restriction (R1), no duplicates, and summed error verified | Not requested; independent acceptance and full downstream audit pending\par
F02 | Major | \(h_*\) satisfies the stated bound’s hypotheses and restores density estimate | Not requested; independent acceptance pending\par
(4.17) | Unrated scope | Unbounded-imaginary uniform reading ruled out; fixed-real use distinguished | Clarify intended constant dependence and range\par
The supported conclusion concerns two invalid derivations in the printed proof. The proposed local corrections preserve the intended statements on the dependencies checked here. No new empirical experiment is required to validate either calculation. External analytic dependencies, a complete independent proof of the final spacing transfer, and independent scientific acceptance remain outside the completed initial review.\par
9 P. Kurlberg and S. Lester, Poisson spacing statistics for lattice points on circles, Invent. Math. 240 (2025), 537–585. Publisher version of record, acquired PDF. https://doi.org/10.1007/s00222-025-01324-1. P. Kurlberg and S. Lester, same title, arXiv:2112.08522v1, 15 December 2021. Targeted released TeX inspection and version metadata. https://arxiv.org/abs/2112.08522.\par
\section*{Supporting checks}
The native source extraction has SHA-256\par
and the related older TeX file has SHA-256\par
The reviewer-written Python check used only standard-library rational and integer arithmetic. It constructed the eight complementary subset-class rows, enumerated their exact rational row-span closures, checked all 18 strata, verified multiplicity 2 versus 1 at \((1,0,3)\), counted 48 points in the stated box, and checked the prime-power hypothesis failure. It did not run author code, simulate Poisson statistics, or replace any analytic proof. Its JSON output SHA-256 is\par
All assertions passed.\par
Native Poppler extraction and page rendering, targeted text searches, the exact source files, and the per-page observations are linked in the JSON provenance. The source PDF pages actually viewed were 1, 3, 4, 7, 9, 12, 14, 15, 23, 24, 33, 34, 41–47. The report is compiled by the authorized portable Tectonic 0.17.0 with `–untrusted`. The compiler executable SHA-256 is\par
The cache snapshot is a superset of used bundle dependencies, not a claim that every cached file was used. Every final report page is rendered and personally inspected before the delivery state is finalized. The saved page-check record binds those observations to the exact PDF.\par
All scientific reasoning, independent calculations, static source inspection, compilation, and image inspection were performed by this assistant. Work’s role is scheduling. The previously confirmed account context is Minhao Liu’s Group–Academic Researchers. No backend equivalence is asserted, and no independent mismatch evidence was found. No browser takeover was attempted in this review.\par
This is initial delivery only: `human\_reviews=[]`, `scientific\_decision=pending`, and independent scientific acceptance is false.\par
\end{document}
